\section{Introduction}
\label{sec:intro}

Public leaderboards shape which language models people download, fine-tune and
deploy.\ifarxiv\else\footnote{The authors used a large language model (Claude, by Anthropic) as an assistant for
writing and debugging analysis code and for drafting and editing the text. Every reported number is
generated by the released code from the released result files, and the authors take full
responsibility for the content.}\fi{} For benchmarks whose answers are free text, such as competition
mathematics, a score comes from two steps. The model writes a response, and then a program called a
\emph{verifier} reads it, finds the answer it gives, and checks that answer against the benchmark's
reference, or \emph{gold}, answer. The leaderboard reports the model's name and its score. It does
not report which verifier produced the score, even though verifiers differ between evaluation tools,
between tasks inside one tool, and over time for a single task.

The choice is not a detail. In February 2025 the Open LLM Leaderboard re-scored its MATH results with
a different verifier \citep{kydlicek2025fixing}. We traced the change file by file through the
leaderboard's own records: it rewrote the scores of runs it had already stored rather than generating
new responses (Appendix~\ref{app:board}). Among the \BoardShiftModels{} models whose scores before and
after come from the same run, \BoardShiftDiscordantPct\% of the pairs swapped order (ties excluded).
The responses were the same, so a ranking that moves this much is partly a ranking of verifiers.

A single response shows how this happens. Asked how many more euros than pounds a person called
Kathy has, one model converted both amounts, compared them, and ended with: ``Kathy has
$\boxed{22}$ more euros than pounds. Final Answer: The final answer is 22. Great job!'' The gold
answer is 22, and a person reading the response would mark it correct. The leaderboard's verifier,
however, looks for the sentence ``Final Answer: The final answer is X. I hope it is correct.'' This
response ended with ``Great job!'' where the verifier expected ``I hope it is correct.'', so the
verifier found no answer and scored the response as wrong. A response in which no answer is found
is recorded exactly like a response with a wrong answer, and the published score does not say which
of the two happened. Whether a model follows the demonstrated format is a legitimate thing to
measure; our question is narrower: what does a ranking say about the answers a reader would find?

Neither verifier is the truth, so a change of verifier shows that a ranking can move without saying
which orderings to trust. A human reader is the reference, but reading every response of every model is
rarely practical, and we are not aware of a measurement of how much of a published ranking a reader would support. We make
that measurement on the archived MATH-Hard board of the Open LLM Leaderboard
\citep{openllmleaderboard, fourrier2024openllm}, which published every model's response to each of
its \NumItems{} competition problems. Three verifiers with different extraction rules read all
\NumRows{} responses of \NumModels{} models. Where all three agree, a human reader almost always
agrees with them; where they disagree, at least one of them is wrong. Most disagreements are not
judgement calls: on \TaxExtraction{} of the \RowsAnyDisagree{} disputed responses, a verifier that gave
no credit had found no answer at all.

We then ask what a human audit of 400 responses can establish about the ranking. For each pair of
models we find every accuracy gap consistent with the verdicts and the labels; if all of them favour
the same model, the ordering is \emph{certified}. The answer, with a guarantee that holds for all
\NumOrderings{} orderings at once and needs no large-sample approximation, comes in two parts. On its
own, the audit certifies \CertFixedNoPooling{} orderings: spread over \NumModels{} models, 400 labels
leave every model's accuracy too loosely pinned down. Letting models share labels where all three
verifiers agree changes this, at the price of an assumption that we state with a single number: the
rate at which a reader overrules three agreeing verifiers may differ between models by at most a
tolerance $\varepsilon$. The audit then certifies \CertFixedPrereg{} orderings at a reference tolerance,\footnote{Our analysis plan is a dated, version-controlled document
released \SuppWhere, not a registration with an external registry, and it was written after
the 400-item audit's labels existed. The reference tolerance is a point it fixes on a curve, not a
confirmatory threshold; Appendix Table~\ref{tab:chronology} orders each choice in the plan against
the labels.} $\varepsilon = \EpsPrereg$, and \CertFixedEpsZero{} at $\varepsilon = 0$. Because each count rests on the assumption,
we report the whole curve over $\varepsilon$ and, for every ordering, a breakdown value: the variation
between models it would take to overturn it.

The paper makes four contributions.
\begin{enumerate}
\item \emph{The verifier moves the whole board.} On unchanged responses, the leaderboard's own change
  of verifier reversed \BoardShiftDiscordantPct\% of the pairwise orderings among
  \BoardShiftModels{} models, and the scoring design behind this is common in four evaluation
  frameworks (Section~\ref{sec:generality}).
\item \emph{Which verifier scored the board is established, not guessed.} Re-running the scoring code
  recorded in its result files reproduces the published verdict on \EmReexecItems{} of
  \EmReexecItems{} responses (Section~\ref{sec:data}), and most disagreements between verifiers are
  failures to find an answer (Section~\ref{sec:disagree}).
\item \emph{A small audit settles nothing on its own.} This audit certifies no ordering
  without an assumption that links models, and every model's rank interval spans all \NumModels{}
  places (Sections~\ref{sec:setup} and~\ref{sec:results}).
\item \emph{Principal result.} Under one stated assumption, how much of a ranking survives a human
  reader can be bounded and reported as a curve, with a breakdown value for each ordering
  (Section~\ref{sec:results}).
\end{enumerate}

Section~\ref{sec:data} describes the data, the verifiers and the human audit;
Section~\ref{sec:disagree} shows what the verifiers disagree about; Section~\ref{sec:setup} develops
the method; Section~\ref{sec:results} gives the answer; Section~\ref{sec:generality} looks beyond the
roster; Sections~\ref{sec:related} and~\ref{sec:limitations} cover related work and limitations. We
make no claim about models outside the roster or about benchmarks other than MATH-Hard.
